% exam2-predicate-logic.tex (slides)

\begin{frame}[allowframebreaks]{测验 2：谓词逻辑}

\begin{problem}[一阶谓词逻辑公式]
  考虑如下欧式空间 $\mathbb{R}^{n}$ 中``开集'' (Open Set) 的定义:

  "A subset $U$ of the Euclidean $n$-space $\mathbb{R}^{n}$
    is {\it open} if, for every point $x$ in $U$,
    there exists a positive real number $\epsilon$
    (depending on $x$) such that any point in $\mathbb{R}^{n}$
    whose Euclidean distance from $x$ is smaller than $\epsilon$
    belongs to $U$."

  \begin{enumerate}[(1)]
    \item 请将"Open Set"（开集）的定义翻译成一阶谓词逻辑公式,
      记为 $\phi$。
      请自行选择论域（全集 Universe）、谓词符号和函数符号, 并请说明它们的含义。
      \begin{solution}
        取论域 $|\mathcal{M}|=\R^{n} \cup \R$。
        \[
          \phi \triangleq
          \forall x \in U \;\exists \epsilon \in \R^{+}\;
            \forall y \in \R^{n}\;
            \bigl(d(y,x) < \epsilon \to y \in U\bigr).
        \]
      \end{solution}
    \item 请给出 $\lnot \phi$ 公式。
      \begin{solution}
        \[
          \begin{aligned}
            \lnot\phi
              \equiv \exists x \in U \;\forall \epsilon \in \R^{+}\;\exists y \in \R^{n}\;
              \bigl(d(y,x) < \epsilon \land y \notin U\bigr).
          \end{aligned}
        \]
      \end{solution}
  \end{enumerate}
\end{problem}

\begin{problem}[一阶谓词逻辑的语法: 替换]
  考虑如下公式 $\phi$:
  \[
    \phi \triangleq \forall x\, ((\exists y\, P(x, y, z)) \land (\forall z\, P(x, y, z))).
  \]
  其中, $P$ 是三元谓词符号。
  请回答以下问题:

  \begin{enumerate}[(a)]
    \item 指出变元 $y$ 的所有自由出现和约束出现。
      \begin{solution}
        将 $y$ 的出现标出:
        \[
          \phi = \forall x\, \bigl((\exists y\, P(x, \blue{y}, z)) \land (\forall z\, P(x, \red{y}, z))\bigr).
        \]
        因此:
        \begin{align*}
          &\text{约束出现: } \exists y\, P(x, \blue{y}, z) \text{ 中的 } y; \\
          &\text{自由出现: } \forall z\, P(x, \red{y}, z) \text{ 中的 } y.
        \end{align*}
      \end{solution}
    \item 令 $t \triangleq g(f(g(y, y)), y)$,
      分别计算 $\phi[t/x]$, $\phi[t/y]$ 与 $\phi[t/z]$。
      \begin{solution}
        记
        \[
          t = g(f(g(y, y)), y).
        \]
        则
        \begin{align*}
          \phi[t/x]
            &= \forall x\, ((\exists y\, P(x, y, z)) \land (\forall z\, P(x, y, z))), \\
          \phi[t/y]
            &= \forall x\, \Bigl((\exists y\, P(x, y, z)) \land
              (\forall z\, P(x, g(f(g(y, y)), y), z))\Bigr), \\
          \phi[t/z]
            &= \forall x\, \Bigl((\exists y\, P(x, y, g(f(g(y, y)), y))) \land
              (\forall z\, P(x, y, z))\Bigr).
        \end{align*}
        其中 $\phi[t/x] = \phi$, 因為 $x$ 在 $\phi$ 中没有自由出现。
      \end{solution}
    \item 分別判断上述 $t$ 在 $\phi$ 中是否可无冲突地替换 $x$、$y$ 与 $z$;
      对于不可无冲突替换的情况, 请说明理由。
      \begin{solution}
        \begin{enumerate}[(i)]
          \item $t$ 对 $x$ 可无冲突替换: $x$ 在 $\phi$ 中无自由出现, 因而不会发生变量捕获;
          \item $t$ 对 $y$ 可无冲突替换: 唯一要替换的自由出现位于 $\forall z\, P(x, y, z)$ 中; 而 $t$ 的自由变元只有 $y$, 不会被 $\forall z$ 捕获。
          \item $t$ 对 $z$ 不可无冲突替换: 要替换的自由出现位于 $\exists y\, P(x, y, z)$ 中; 但 $t$ 含有自由变元 $y$, 替换后会落入 $\exists y$ 的作用域, 从而使 $t$ 中的 $y$ 被捕获。
        \end{enumerate}
      \end{solution}
  \end{enumerate}
\end{problem}

\begin{problem}[一阶谓词逻辑公式的语义]
  考虑公式 $\phi \triangleq \forall x\, \forall y\, \exists z\, (R(x, y) \to R(y, z))$,
  其中 $R$ 是二元谓词符号。
  请判断下列结构 $\mathcal{M}$ 是否满足 $\phi$,
  并说明理由:
  \begin{enumerate}[(1)]
    \item 令 $|\mathcal{M}| = \{a, b, c, d\}$ 且
      $R^{\mathcal{M}} = \{(b, c), (b, b), (b, a)\}$。
      \begin{solution}
        取
        \[
          x = b, \qquad y = c.
        \]
        则
        \[
          R^{\mathcal{M}}(b, c) \text{ 为真}.
        \]
        因而
        \[
          \exists z\, (R(b, c) \to R(c, z))
        \]
        等价于
        \[
          \exists z\, R(c, z).
        \]
        但由
        \[
          R^{\mathcal{M}} = \{(b, c), (b, b), (b, a)\}
        \]
        可知不存在任何 $z \in |\mathcal{M}|$ 使得 $(c, z) \in R^{\mathcal{M}}$。
        故
        \[
          \mathcal{M} \not\models \phi.
        \]
      \end{solution}
    \item 令 $|\mathcal{M}| = \{a, b, c\}$ 且
      $R^{\mathcal{M}} = \{(b, c), (a, b), (c, b)\}$。
      \begin{solution}
        只需检查所有使 $R(x, y)$ 为真的二元组;
        对每个这样的 $y$, 给出一个见证 $z$ 即可。

        需要检查的真元组只有
        \[
          (a, b), \qquad (b, c), \qquad (c, b).
        \]
        分別取见证:
        \begin{align*}
          &(x, y) = (a, b): && z = c, && R(b, c) \text{ 为真}; \\
          &(x, y) = (b, c): && z = b, && R(c, b) \text{ 为真}; \\
          &(x, y) = (c, b): && z = c, && R(b, c) \text{ 为真}.
        \end{align*}
        故
        \[
          \mathcal{M} \models \phi.
        \]
      \end{solution}
  \end{enumerate}
\end{problem}

\begin{problem}[一阶谓词逻辑的自然推理系统]
  请使用自然推理规则证明:

  \begin{enumerate}[(1)]
    \setcounter{enumi}{2}
    \item
      \[
        \exists x\, (P(x) \land Q(x)) \vdash \exists x\, P(x) \land \exists x\, Q(x)
      \]
  \end{enumerate}
\end{problem}

\begin{logicproof}{4}
  \exists x\, (P(x) \land Q(x)) & premise \\
  \begin{subproof}
    x_0\; (\text{\blue{fresh}}),\; P(x_0) \land Q(x_0) & assumption \\
    P(x_0) & $\land\elimleft$ 2 \\
    Q(x_0) & $\land\elimright$ 2 \\
    \exists x\, P(x) & $\exists$-intro 3 \\
    \exists x\, Q(x) & $\exists$-intro 4 \\
    \exists x\, P(x) \land \exists x\, Q(x) & $\land\intro$ 5, 6
  \end{subproof}
  \exists x\, P(x) \land \exists x\, Q(x) & $\exists$-elim 1, 2--7
\end{logicproof}

\end{frame}
